\documentclass[11pt,a4paper]{article}
\usepackage[utf8]{inputenc}
\usepackage[T1]{fontenc}
\usepackage[spanish]{babel}
\usepackage{amsmath,amssymb,mathtools}
\usepackage{geometry}
\usepackage{enumitem}
\usepackage{lmodern}
\usepackage{microtype}
\geometry{margin=2.4cm}
\setlength{\parindent}{0pt}
\setlength{\parskip}{0.45em}
\newcommand{\R}{\mathbb{R}}
\newcommand{\C}{\mathbb{C}}
\newcommand{\K}{\mathbb{K}}
\newcommand{\Ker}{\operatorname{Ker}}
\newcommand{\rg}{\operatorname{rg}}
\newcommand{\sg}{\operatorname{sg}}
\newcommand{\Span}{\operatorname{L}}
\begin{document}

\begin{center}
{\LARGE\bfseries Álgebra Lineal}\\[0.4em]
{\Large Tema 5: Formas Bilineales y Cuadráticas}
\end{center}

\tableofcontents
\newpage

\section{Aplicaciones multilineales y formas bilineales}

Sean $V_1,\ldots,V_n,W$ espacios vectoriales sobre $\K$.

Una aplicación:

\[
f:V_1\times\cdots\times V_n\longrightarrow W
\]

es multilineal si es lineal en cada una de sus componentes.

Si:

\[
f:V\times V\longrightarrow\K,
\]

se denomina \textbf{forma bilineal} cuando es lineal en cada variable:

\[
f(au+bv,w)=af(u,w)+bf(v,w)
\]

\[
f(u,av+bw)=af(u,v)+bf(u,w).
\]

\subsection{Propiedades}

Toda forma bilineal verifica:

\[
f(u,0)=f(0,v)=0
\]

\[
f(-u,v)=f(u,-v)=-f(u,v)
\]

y:

\[
f\left(\sum_i a_i u_i,\sum_j b_jv_j\right)
=
\sum_{i,j}a_ib_jf(u_i,v_j).
\]

\section{Formas bilineales simétricas y antisimétricas}

Una forma bilineal es \textbf{simétrica} si:

\[
\boxed{
f(x,y)=f(y,x)
}
\]

y es \textbf{antisimétrica} si:

\[
\boxed{
f(x,y)=-f(y,x).
}
\]

Sobre $\R$ o $\C$:

\[
\boxed{
f\text{ antisimétrica}
\iff
f(x,x)=0\quad\forall x.
}
\]

Toda forma bilineal puede descomponerse como:

\[
\boxed{
f=f_s+f_a
}
\]

donde:

\[
\boxed{
f_s(u,v)=\frac12\left(f(u,v)+f(v,u)\right)
}
\]

y:

\[
\boxed{
f_a(u,v)=\frac12\left(f(u,v)-f(v,u)\right).
}
\]

\section{Matriz asociada a una forma bilineal}

Sea:

\[
B=\{v_1,\ldots,v_n\}
\]

una base de $V$.

La matriz asociada a $f$ en $B$ es:

\[
\boxed{
M_B(f)=\left(f(v_i,v_j)\right).
}
\]

Si:

\[
[x]_B=X,
\qquad
[y]_B=Y,
\]

entonces:

\[
\boxed{
f(x,y)=X^tM_B(f)Y.
}
\]

En particular, una matriz cuadrada $A$ define la forma bilineal:

\[
\boxed{
f(x,y)=X^tAY.
}
\]

Además:

\[
f\text{ simétrica}
\iff
A^t=A
\]

y:

\[
f\text{ antisimétrica}
\iff
A^t=-A.
\]

\section{Cambio de base y congruencia}

Sean $B$ y $B'$ dos bases y:

\[
X=PX',
\qquad
Y=PY'.
\]

Entonces:

\[
f(x,y)
=
X'^tP^tAPY'.
\]

Por tanto:

\[
\boxed{
M_{B'}(f)=P^tM_B(f)P.
}
\]

Dos matrices que representan la misma forma bilineal en bases distintas son \textbf{congruentes}.

Las matrices congruentes tienen el mismo rango.

\section{Rango, núcleo y degeneración}

El rango de una forma bilineal es:

\[
\boxed{
\rg f=\rg M_B(f).
}
\]

Su núcleo es:

\[
\boxed{
N(f)
=
\{x\in V:f(x,y)=0\quad\forall y\in V\}.
}
\]

Matricialmente:

\[
\boxed{
N(f)=\Ker M_B(f).
}
\]

Si:

\[
\dim V=n,
\]

la forma es \textbf{no degenerada} si:

\[
\boxed{
\rg f=n
}
\]

y es \textbf{degenerada} si:

\[
\boxed{
\rg f<n.
}
\]

\section{Formas cuadráticas}

Si $f:V\times V\to\K$ es bilineal, la aplicación:

\[
q:V\longrightarrow\K
\]

definida por:

\[
\boxed{
q(v)=f(v,v)
}
\]

es una forma cuadrática.

Una misma forma cuadrática puede proceder de varias formas bilineales, pero existe una única forma bilineal simétrica asociada.

\section{Forma polar}

La forma polar asociada a una forma cuadrática $q$ es:

\[
\boxed{
f_q(u,v)
=
\frac12
\left(
q(u+v)-q(u)-q(v)
\right).
}
\]

Es bilineal y simétrica, y satisface:

\[
\boxed{
f_q(v,v)=q(v).
}
\]

Si una forma bilineal $f$ genera $q$, su forma polar es precisamente la parte simétrica de $f$.

\section{Matriz asociada a una forma cuadrática}

La matriz asociada a una forma cuadrática es la matriz de su forma polar.

Si:

\[
A=M_B(q),
\]

entonces $A$ es simétrica y:

\[
\boxed{
q(x)=X^tAX.
}
\]

Para construirla a partir de la expresión de $q$:

\begin{itemize}
\item los coeficientes de $x_i^2$ se colocan en la diagonal;
\item el coeficiente de $x_ix_j$ se divide entre dos y se coloca en las posiciones $(i,j)$ y $(j,i)$.
\end{itemize}

\section{Conjugación}

Sea $f$ una forma bilineal simétrica.

Dos vectores $u,v$ son \textbf{conjugados} si:

\[
\boxed{
f(u,v)=0.
}
\]

Un vector no nulo es autoconjugado si:

\[
f(v,v)=0.
\]

El núcleo es el conjunto de vectores conjugados con todos los vectores del espacio.

Para un subconjunto no vacío $S\subseteq V$, su conjugado es:

\[
\boxed{
S^c
=
\{u\in V:f(u,v)=0\quad\forall v\in S\}.
}
\]

Si:

\[
S=\Span\{v_1,\ldots,v_k\},
\]

entonces:

\[
x\in S^c
\iff
f(x,v_i)=0
\quad
i=1,\ldots,k.
\]

Se verifica:

\[
\boxed{
\dim S^c+\dim S\geq n.
}
\]

Si la forma es no degenerada:

\[
\boxed{
\dim S^c+\dim S=n.
}
\]

\section{Diagonalización por congruencia}

Toda forma bilineal simétrica posee una base de vectores conjugados.

Por tanto, toda matriz simétrica es congruente con una matriz diagonal:

\[
\boxed{
D=P^tAP.
}
\]

Las columnas de $P$ son las coordenadas de una base de vectores conjugados.

En dicha base:

\[
\boxed{
q(x_1,\ldots,x_n)
=
d_1x_1^2+\cdots+d_nx_n^2.
}
\]

Sobre $\R$, mediante una normalización adecuada, toda forma cuadrática es congruente con una diagonal cuyos elementos son:

\[
\boxed{
1,\;-1,\;0.
}
\]

\subsection{Método por transformaciones elementales}

Para diagonalizar una matriz simétrica $A$:

\begin{enumerate}
\item se aplican operaciones elementales por filas para triangularizar;
\item se aplican las mismas operaciones por columnas;
\item las operaciones aplicadas a $I$ permiten obtener la matriz de paso.
\end{enumerate}

El resultado cumple:

\[
\boxed{
D=P^tAP.
}
\]

\section{Ley de inercia de Sylvester}

Todas las matrices diagonales asociadas a una misma forma cuadrática real tienen el mismo número de elementos positivos, negativos y nulos.

Si hay:

\[
p\text{ positivos},
\qquad
q\text{ negativos},
\]

la signatura es:

\[
\boxed{
\sg(q)=(p,q).
}
\]

Además:

\[
\boxed{
\rg q=p+q.
}
\]

\section{Clasificación de formas cuadráticas reales}

Una forma cuadrática real es:

\[
\boxed{
\text{definida positiva}
\iff
\sg(q)=(n,0)
}
\]

\[
\boxed{
\text{semidefinida positiva}
\iff
\sg(q)=(p,0),\quad p<n
}
\]

\[
\boxed{
\text{definida negativa}
\iff
\sg(q)=(0,n)
}
\]

\[
\boxed{
\text{semidefinida negativa}
\iff
\sg(q)=(0,q),\quad q<n
}
\]

\[
\boxed{
\text{indefinida}
\iff
p>0,\ q>0.
}
\]

\section{Criterio de Sylvester}

Sean:

\[
\Delta_k
\]

los menores principales de orden $k$ de una matriz simétrica $A$.

La forma es definida positiva si:

\[
\boxed{
\Delta_k>0
\quad
\forall k=1,\ldots,n.
}
\]

Es definida negativa si:

\[
\boxed{
(-1)^k\Delta_k>0
\quad
\forall k=1,\ldots,n.
}
\]

En dimensión $3$, si:

\[
\Delta_1>0,
\qquad
\Delta_2>0,
\qquad
\Delta_3=0,
\]

la forma es semidefinida positiva.

\section{Determinación práctica del carácter}

El carácter de una forma cuadrática puede determinarse mediante:

\begin{enumerate}
\item una diagonalización por congruencia;
\item los menores principales;
\item el signo de los autovalores de una matriz asociada;
\item el rango y la signatura.
\end{enumerate}

\section{Desigualdad de Schwarz}

Si $q$ es definida positiva y $f_q$ es su forma polar:

\[
\boxed{
f_q(u,v)^2
\leq
q(u)q(v).
}
\]

\section{Formas sesquilineales}

Sea $V$ un espacio vectorial complejo.

Una forma sesquilineal:

\[
f:V\times V\longrightarrow\C
\]

es lineal en la primera variable y semilineal en la segunda:

\[
f(au,v)=af(u,v)
\]

\[
f(u,bv)=\overline b\,f(u,v).
\]

Si:

\[
[x]_B=X,
\qquad
[y]_B=Y,
\]

entonces:

\[
\boxed{
f(x,y)=X^tA\overline Y.
}
\]

Bajo un cambio de base $X=PX'$:

\[
\boxed{
A'=P^tA\overline P.
}
\]

Una forma sesquilineal es hermítica si:

\[
\boxed{
f(v,u)=\overline{f(u,v)}.
}
\]

Equivalentemente, su matriz es hermítica:

\[
\boxed{
\overline A^t=A.
}
\]

\end{document}
